\section{What a video JEPA learns from innate-immune video}
\label{sec:partA}

\subsection{Morphology of activated microglia}
On BBBC054, every learned representation separates the three microglia morphology classes
far better than hand-crafted shape and texture features (\val{bb.E1_handcrafted_logreg.3c})
or exposure time alone (\val{bb.time.3c}; Figure~\ref{fig:benchmarks}). The pretrained
video model V-JEPA~2 is the best encoder (\val{bb.E7_vjepa2_vitl.3c}). TiJepa-1
(\val{bb.E4_tijepa1.3c}) is not distinguishable from DINOv2 (\val{bb.E3_dinov2_vit_s.3c}),
the ImageNet ViT (\val{bb.E2_imagenet_vit_s.3c}) or TiJepa-0 (\val{bb.E0_tijepa0.3c});
V-JEPA~2 exceeds it by \val{bb.d.E4_tijepa1.E7_vjepa2_vitl.3c.abs} (95\% CI
\val{bb.ci.E4_tijepa1.E7_vjepa2_vitl.3c.abs}). Training on cell video therefore did not beat a
general-purpose video model, and twice as many training steps on a quality-filtered corpus
did not improve on TiJepa-0. Within the TiJepa recipe, future masking adds
\val{bb.d.E4_tijepa1.E4s_tijepa1_spatial.3c}
(\val{bb.ci.E4_tijepa1.E4s_tijepa1_spatial.3c}) over the spatial-only ablation and training
adds \val{bb.d.E4_tijepa1.E6_random_init.3c} (\val{bb.ci.E4_tijepa1.E6_random_init.3c})
over a randomly initialised network of the same architecture, which itself reaches
\val{bb.E6_random_init.3c}. In the round-versus-amoeboid task the ordering is the same at a
higher level (V-JEPA~2 \val{bb.E7_vjepa2_vitl.ra}, TiJepa-1 \val{bb.E4_tijepa1.ra}, DINOv2
\val{bb.E3_dinov2_vit_s.ra}). Accuracy depends on apparent cell size: rescaling the crops
by 0.5, 0.75, 1 and 1.5 before encoding gives \val{bb.scale.0.5}, \val{bb.scale.0.75},
\val{bb.scale.1.0} and \val{bb.scale.1.5} for TiJepa-1, so differences in magnification
between pretraining and evaluation data matter.

\subsection{Label-free cell identity}
On T cells and neutrophils imaged in the same fields, frozen encoders reach balanced
accuracies of \val{af.Am.frozen.min}--\val{af.Am.frozen.max} on the \val{af.Am.fovs}
mixed fields, every one of them inside the 95\% CI of a probe on per-channel intensity
statistics (\val{af.Am.intensity_baseline.all}, CI \val{af.Am.intensity_baseline.all.ci};
Figure~\ref{fig:benchmarks}). TiJepa-1 reaches
\val{af.Am.E4_tijepa1.all} and V-JEPA~2 \val{af.Am.E7_vjepa2_vitl.all}; random
initialisation reaches \val{af.Am.E6_random_init_draw0.all}. On the gradient-contrast
channel alone, which is closest to the brightfield TiJepa was trained on, TiJepa-1 gives
\val{af.Am.E4_tijepa1.dodt} and V-JEPA~2 \val{af.Am.E7_vjepa2_vitl.dodt}. Cell identity in
these data is carried by autofluorescence intensity; frozen image features add nothing measurable
to it. In the six-class set, where each cell type was imaged in its own session, the
hand-crafted features reach \val{af.B.E1_handcrafted.all} and TiJepa-1
\val{af.B.E4_tijepa1.all}; that comparison cannot separate cell type from session.

\subsection{NAD(P)H lifetime from the photon image}
In \val{fl.cells} macrophages from \val{fl.fovs} fields, the per-cell mean NAD(P)H lifetime
(\val{fl.tau.mean}\,$\pm$\,\val{fl.tau.sd}\,ns; our data check required a median between
0.3 and 3.5\,ns)
is predicted from frozen features of the photon-count image with $R^2$ up to
\val{fl.tau.E4s_tijepa1_spatial} (TiJepa-1 \val{fl.tau.E4_tijepa1}, V-JEPA~2
\val{fl.tau.E7_vjepa2_vitl}), against \val{fl.tau.intensity_baseline} for intensity statistics.
That gain does not reflect cell-level metabolism. Lifetime differs between experimental
arms, and a reference that predicts each cell by its arm's training-fold mean reaches
$R^2=\val{fl.ref.tau}$, more than any image feature. When targets and features are centred
on their field's mean, so that only variation between cells of one field counts, TiJepa-1
reaches \val{fl.wf.tau.E4_tijepa1} (95\% CI \val{fl.wf.tau.E4_tijepa1.ci}), V-JEPA~2
\val{fl.wf.tau.E7_vjepa2_vitl} and intensity statistics \val{fl.wf.tau.intensity_baseline}
(\val{fl.wf.tau.intensity_baseline.ci}); the bound fraction behaves the same way. These two
controls were added after seeing the first result. The frozen features recognise the
condition a field came from; they do not read a cell's lifetime beyond its brightness.

\subsection{Forecasting is explained by positional structure}
\label{sec:forecasting}
TiJepa-1's predictor forecasts its own future latents with an error \val{fc.E4_tijepa1.1.ratio_vs_copy_last}
times that of copying the last context latent, against \val{fc.E4s_tijepa1_spatial.1.ratio_vs_copy_last}
for the spatial-only ablation, which never saw a future mask, and the advantage holds at
every stride (Figure~\ref{fig:assoc}). It also beats a slot-conditional mean fitted on other
movies (ratio \val{fc.E4_tijepa1.1.ratio_vs_slot_mean}, exact $p=\val{fc.E4_tijepa1.1.p.slot}$,
the smallest possible with six movies), so its predictions are specific to the clip. But
the controls show that the advantage over copy-last is not evidence of modelling cell
dynamics. On static clips, in which nothing moves, the predictor keeps
\val{fc.E4_tijepa1.1.adv.static} of its \val{fc.E4_tijepa1.1.adv} improvement over
copy-last; given the context of a different clip from the same movie it keeps
\val{fc.E4_tijepa1.1.adv.swap}. Copy-last is a weak reference because the target encoder's
latents change with temporal position even when the image does not, and the predictor
receives those positions. Decoded to pixels, the predictions are no closer to the future
frames than the last observed frame (error ratio \val{fc.E4_tijepa1.1.pixel_ratio_decoded_pred_vs_pixel_copy_last}
at stride 1, \val{fc.E4_tijepa1.27.pixel_ratio_decoded_pred_vs_pixel_copy_last} at stride
27). Our pre-specified criterion for forecasting \val{fc.pass}, and we make no forecasting
claim.
